\import{test/phase7/composition-forms-base.tex}

\begin{proposition}\label{calculation_form_parameters}
    Suppose $a\calceq{n}b\calceq{n}c\calceq{n}d$.
    Then $a\calceq{n}d$.
\end{proposition}
\begin{proof}
    \begin{align*}
        a
            &\calceq{n} b
                \explanation{by assumption}
                \\
            &\calceq{n} c
                \explanation{by assumption}
                \\
            &\calceq{n} d
                \explanation{by assumption}
    \end{align*}
\end{proof}

\begin{proposition}\label{calculation_form_coinciding_capture}
    Suppose $a=b$ and $b\calceq{a}c$.
    Then $a\calceq{a}c$.
\end{proposition}
\begin{proof}
    \begin{align*}
        a
            &= b
                \explanation{by assumption}
                \\
            &\calceq{a} c
                \explanation{by assumption}
    \end{align*}
\end{proof}

\begin{proposition}\label{calculation_form_negative}
    Suppose $x\notin B$ and $B\supseteq A$.
    Then $x\notin A$.
\end{proposition}
\begin{proof}
    \begin{align*}
        x
            &\notin B
                \explanation{by assumption}
                \\
            &\supseteq A
                \explanation{by assumption}
    \end{align*}
\end{proof}

\begin{proposition}\label{calculation_form_disequality_left}
    Suppose $a=b$ and $b\neq c$.
    Then $a\neq c$.
\end{proposition}
\begin{proof}
    \begin{align*}
        a
            &= b
                \explanation{by assumption}
                \\
            &\neq c
                \explanation{by assumption}
    \end{align*}
\end{proof}

\begin{proposition}\label{calculation_form_disequality_right}
    Suppose $a\neq b$ and $b=c$.
    Then $a\neq c$.
\end{proposition}
\begin{proof}
    \begin{align*}
        a
            &\neq b
                \explanation{by assumption}
                \\
            &= c
                \explanation{by assumption}
    \end{align*}
\end{proof}

\begin{proposition}\label{calculation_form_reversed}
    Suppose $a=b$ and $b\calcreverse c$.
    Then $a\calcreverse c$.
\end{proposition}
\begin{proof}
    \begin{align*}
        a
            &= b
                \explanation{by assumption}
                \\
            &\calcreverse c
                \explanation{by assumption}
    \end{align*}
\end{proof}

\begin{proposition}\label{calculation_form_ignored_endpoint}
    Suppose $b=c$.
    Then $a\calcignore c$.
\end{proposition}
\begin{proof}
    \begin{align*}
        a
            &\calcignore b
                \\
            &= c
                \explanation{by assumption}
    \end{align*}
\end{proof}

\begin{proposition}\label{calculation_form_carrier}
    Let $G,H$ be quasiordered sets. Suppose $a\in G$ and $G=H$.
    Then $a\in H$.
\end{proposition}
\begin{proof}
    \begin{align*}
        a
            &\in G
                \explanation{by assumption}
                \\
            &= H
                \explanation{by assumption}
    \end{align*}
\end{proof}

\begin{proposition}\label{calculation_form_negative_carrier}
    Let $G,H$ be quasiordered sets. Suppose $a\notin G$ and $G=H$.
    Then $a\notin H$.
\end{proposition}
\begin{proof}
    \begin{align*}
        a
            &\notin G
                \explanation{by assumption}
                \\
            &= H
                \explanation{by assumption}
    \end{align*}
\end{proof}

\begin{proposition}\label{calculation_form_conditions}
    Suppose $a,b,c\in A$. Suppose $a\calcguard{A}b\calcguard{A}c$.
    Then $a\calcguard{A}c$.
\end{proposition}
\begin{proof}
    \begin{align*}
        a
            &\calcguard{A} b
                \explanation{by assumption}
                \\
            &\calcguard{A} c
                \explanation{by assumption}
    \end{align*}
\end{proof}

\begin{proposition}\label{calculation_form_notion_publication}
    Suppose $a\in A$ and $A=B$.
    Then $a\in B$.
\end{proposition}
\begin{proof}
    \begin{align*}
        a
            &\in A
                \explanation{by notion}
                \\
            &= B
                \explanation{by assumption}
    \end{align*}
    Follows by notion.
\end{proof}

\begin{proposition}\label{calculation_form_term_single}
    Suppose $a\mathrel{R}b$.
    Then $a\mathrel{R}b$.
\end{proposition}
\begin{proof}
    \begin{align*}
        a
            &\mathrel{R} b
                \explanation{by assumption}
    \end{align*}
\end{proof}

\begin{proposition}\label{calculation_prefix_formation}
    Let $A$ be a set.
    For all $x,y\in A$ such that $x=y$ we have $x\calctyped{A}y$.
\end{proposition}
\begin{proof}
    For all $x,y\in A$ such that $x=y$, we have
    \begin{align*}
        x
            &\calctyped{A}y
                \explanation{by \obv{calculation_typed_relation}}
    \end{align*}
    Follows.
\end{proof}
